\section{Marco Contextual}

El marco contextual delimita el entorno geográfico, económico, tecnológico e institucional en el cual se inserta la problemática investigada, permitiendo caracterizar tanto las tendencias observadas a nivel global como las particularidades operativas del escenario local en el que se proyecta la propuesta (Hernández Sampieri, 2014).

\subsection{Contexto internacional del soporte técnico y la trazabilidad de servicios TI}

En el ámbito internacional de las economías industrializadas, la prestación de servicios de soporte técnico se encuentra altamente automatizada mediante plataformas corporativas de Enterprise Resource Planning (ERP) y módulos especializados de gestión de servicios, tales como SAP Service Cloud (SAP SE, 2023). Asimismo, en el segmento de pequeñas y medianas empresas se ha masificado el uso de soluciones modulares e integradas de código abierto o modelo freemium/SaaS como Odoo, la cual incluye módulos de mesa de ayuda (Helpdesk), inventarios y reparación (Odoo S.A., 2024). Estas herramientas demuestran que la trazabilidad integral incrementa la satisfacción del usuario y reduce los costos operativos de búsqueda documental. No obstante, el análisis de sus funcionalidades evidencia que plataformas como SAP o los módulos avanzados de Odoo requieren costos de licenciamiento en divisas, procesos complejos de configuración y administración continua por personal de TI dedicado, características que las vuelven poco idóneas o costosas para talleres independientes de reparación de computadoras en entornos emergentes.

\subsection{Contexto latinoamericano y brechas operativas en micro y pequeñas empresas}

En el contexto latinoamericano, si bien las empresas medianas y grandes han adoptado plataformas corporativas de soporte e infraestructura administrada a través de proveedores regionales como Sonda o Noventiq (Sonda S.A., 2024; Noventiq, 2024), existe una marcada brecha tecnológica en las micro y pequeñas empresas (MIPYMES) de mantenimiento informático. Estudios como el de Mosquera Perea et al. (2025) en Colombia evidencian que los talleres independientes de reparación sufren de desorganización en el registro de recepción y falta de mecanismos para demostrar la transparencia en las cotizaciones. Asimismo, soluciones de código abierto populares en la región como GLPI (GLPI Project, 2024), si bien destacan en la gestión de inventarios y tickets de TI corporativos internos, carecen de canales conversacionales con inteligencia artificial orientados al cliente final y de flujos ágiles de cotización y evidencias fotográficas para talleres comerciales independientes.

\subsection{Contexto nacional (Bolivia) y madurez del ecosistema digital}

En Bolivia, el ecosistema de base digital formal presenta una marcada concentración geográfica en las ciudades del eje troncal: Santa Cruz (40,9\%), La Paz (33,8\%) y Cochabamba (20,1\%), acumulando el 94,8\% del total nacional de emprendimientos digitales registradas en el reporte del Mapeo del Ecosistema de Tecnología Digital en Bolivia 2026 (Fundación Emprender Futuro et al., 2026). A nivel nacional existen instituciones promotoras de la trazabilidad empresarial como GS1 Bolivia (GS1 Bolivia, 2026), la cual brinda estándares de identificación y código de barras. Asimismo, en la gestión de servicios informáticos destacan soluciones Open Source como GLPI (GLPI Project, 2024) y software de gestión especializado como ExpertoSis, Freshservice o Jira Service Management. No obstante, estas plataformas están orientadas a inventarios institucionales o mesas de ayuda corporativas y no a la dinámica comercial directa entre un taller técnico independiente y sus clientes particulares.

\subsection{Contexto local (Ciudad de Tarija) y dinámicas operativas de los talleres independientes}

En el radio urbano de la ciudad de Tarija, la presente investigación se orienta al estudio de los prestadores de servicios de soporte técnico y mantenimiento de computadoras, abarcando tanto talleres comerciales establecidos como iniciativas y técnicos independientes. Como parte del diagnóstico de campo planificado, se indagará de manera directa sobre las modalidades operativas utilizadas para registrar la recepción de equipos, documentar los diagnósticos, comunicar los presupuestos y dar seguimiento a las reparaciones.

En este marco, se examinará empíricamente el grado de adopción de registros físicos en papel, herramientas digitales locales o aplicaciones de mensajería informal, evitando asumir anticipadamente la prevalencia de una modalidad específica. Asimismo, la ausencia de un directorio gremial o registro comercial centralizado se considera una condición del entorno que justifica la necesidad de realizar un diagnóstico directo. Los hallazgos del trabajo de campo permitirán fundamentar la definición de requerimientos del sistema propuesto en las dinámicas reales identificadas en el contexto tarijeño.

A continuación, se presenta un cuadro comparativo entre las principales soluciones tecnológicas de gestión de servicios y la propuesta de diseño formulada para la ciudad de Tarija:

\begin{table}[H]
\centering
\footnotesize
\setlength{\tabcolsep}{3pt}
\renewcommand{\arraystretch}{1.15}
\begin{tabular}{|p{2.3cm}|p{2cm}|p{2cm}|p{2.9cm}|p{2.5cm}|}
\hline
\textbf{Solución / Plataforma} & \textbf{Licencia / Modelo} & \textbf{Dominio Principal} & \textbf{Trazabilidad e IA} & \textbf{Adecuación a Talleres Tarija} \\
\hline
Odoo (Helpdesk/Repair) & Freemium / SaaS Modulares & ERP General para PYMES & Trazabilidad de tickets; IA mediante módulos pagados. & Parcial: Exige configuración compleja de múltiples módulos. \\
\hline
SAP Service Cloud & De Paga (Corporativa) & Gran Empresa y Corporaciones & Trazabilidad avanzada; IA predictiva empresarial. & Baja: Costos de licenciamiento inaccesibles para talleres pyme. \\
\hline
GLPI (Open Source) & Código Abierto (GPLv3) & Gestión TI e Inventarios de Red & Trazabilidad por tickets; IA vía plugins externos. & Media: Enfoque en departamentos TI internos, no en cliente final. \\
\hline
GS1 Bolivia & Estándares Internacionales & Trazabilidad Logística y Código de Barras & Estándares de identificación y códigos GS1; sin enfoque en tickets o IA. & Baja: Orientado a cadenas de suministro e industria general, no a talleres informáticos. \\
\hline
ExpertoSis (Sistemas de Gestión TI) & Comercial (Propietario) & Gestión de Incidencias y Servicios TI & Gestión de tickets y módulo de trazabilidad (según especificaciones comerciales); sin IA conversacional. & Media: Enfocado en soporte empresarial genérico, sin asistentes conversacionales desacoplados. \\
\hline
Sistema Propuesto (Tarija) & Licencia Propia (Diseño a Medida) & Soporte Técnico de Computadoras & Diseño de trazabilidad por estados secuenciales y evidencias fotográficas; propuesta de asistente de IA desacoplado para consultas de estado e información. & Alta: Diseñado específicamente para las dinámicas de los talleres independientes de Tarija. \\
\hline
\end{tabular}
\caption{Cuadro comparativo de soluciones informáticas de gestión y trazabilidad de soporte técnico.}
\label{tab:comparativo-soluciones}
\end{table}